% Bagian: sets-functions-relations
% Bab: arithmetization
% Subbagian: potongan
%
\documentclass[../../../include/open-logic-section]{subfiles}

\begin{document}

\olfileid[id]{sfr}{arith}{cuts}
\olsection{Dari $\Rat$ ke $\Real$}

Pada dasarnya, Sifat Kelengkapan menunjukkan bahwa setiap titik $\alpha$ pada
garis bilangan real membagi garis itu secara sempurna menjadi dua bagian:
bagian yang batas atas terkecilnya adalah $\alpha$, dan bagian yang batas bawah
terbesarnya adalah $\alpha$. Untuk \emph{membangun} bilangan real dari
bilangan rasional, Dedekind mengusulkan agar kita cukup memandang bilangan
real sebagai \emph{potongan} yang membagi bilangan rasional. Dengan kata lain,
kita mengidentifikasi $\sqrt{2}$ dengan \emph{potongan} yang memisahkan
bilangan rasional $< \sqrt{2}$ dari bilangan rasional $> \sqrt{2}$.

Mari kita rapikan gagasan ini. Jika kita memotong bilangan rasional menjadi
dua bagian, kita dapat mengidentifikasi secara unik partisi yang kita buat
hanya dengan memperhatikan bagian \emph{bawahnya}. Maka, secara saksama, kita
memberikan definisi berikut:

\begin{defn}[Potongan]
Suatu \emph{potongan} $\alpha$ adalah setiap segmen awal sejati takkosong dari
bilangan rasional yang tidak mempunyai elemen terbesar.
Dengan kata lain, $\alpha$ merupakan suatu potongan jika dan hanya jika:
\begin{enumerate}
	\item \emph{takkosong, sejati}: $\emptyset \neq \alpha \subsetneq \Rat$
	\item \emph{awal}: untuk semua $p,q \in \Rat$: jika $p < q \in \alpha$ maka $p \in \alpha$
	\item \emph{tanpa maksimum}: untuk setiap $p \in \alpha$ terdapat $q \in \alpha$ sedemikian sehingga $p < q$
\end{enumerate}
Kemudian $\Real$ adalah himpunan semua potongan.
\end{defn}

Sekarang kita dapat mengatakan bahwa $\sqrt{2} = \Setabs{p \in \Rat}{p^2 < 2\text{
atau }p < 0}$. Tentu saja, kita perlu memeriksa bahwa objek ini \emph{memang}
suatu potongan, tetapi hal itu kita tempatkan dalam \olref[check]{sec}.

Seperti sebelumnya, setelah mendefinisikan beberapa entitas, selanjutnya kita
perlu mendefinisikan fungsi dan relasi dasar padanya. Kita mulai dengan yang
mudah:
\begin{align*}
	\alpha \leq \beta \text{ jika dan hanya jika }\alpha \subseteq \beta
\end{align*}
Definisi urutan ini memungkinkan kita \emph{menyatakan} hasil utama, yaitu
bahwa himpunan semua potongan memiliki Sifat Kelengkapan. Jika dijabarkan
sepenuhnya, pernyataannya berbentuk sebagai berikut. Jika $S$ adalah himpunan
takkosong yang terdiri atas potongan-potongan dan mempunyai suatu batas atas, maka $S$ mempunyai
suatu batas atas terkecil. Secara lebih rinci: terdapat suatu potongan,
$\lambda$, yang merupakan batas atas bagi $S$, yakni\ $(\forall \alpha \in S)\alpha \subseteq
\lambda$, dan $\lambda$ adalah potongan terkecil semacam itu, yakni\ $(\forall \beta \in \Real)((\forall \alpha \in S)\alpha \subseteq \beta \lif \lambda \subseteq \beta)$. Berikut
pembuktian hasil tersebut:

\begin{thm}\ollabel{realcompleteness}
Himpunan semua potongan memiliki Sifat Kelengkapan.
\end{thm}

\begin{proof}
Misalkan $S$ adalah sembarang himpunan takkosong yang terdiri atas potongan-potongan dan mempunyai
suatu batas atas. Misalkan $\lambda = \bigcup S$.
%\Setabs{p \in \Rat}{(\exists \alpha \in S)p \in \alpha}$$
Pertama-tama kita menyatakan bahwa $\lambda$ adalah suatu potongan:
\begin{enumerate}
\item Karena $S$ takkosong, sekurang-kurangnya satu potongan
termasuk dalam $S$, sehingga $\emptyset \neq \lambda$. Karena $S$ adalah
himpunan yang terdiri atas potongan-potongan, $\lambda \subseteq \Rat$. Karena $S$ mempunyai suatu
batas atas, terdapat suatu $p \in \Rat$ yang tidak termasuk dalam satu pun
potongan $\alpha \in S$. Jadi $p\notin \lambda$, sehingga $\lambda \subsetneq
\Rat$.
\item Andaikan $p < q \in \lambda$. Maka terdapat suatu $\alpha \in S$
sedemikian sehingga $q \in \alpha$. Karena $\alpha$ adalah suatu potongan,
$p \in \alpha$. Jadi $p \in \lambda$.
\item Andaikan $p \in \lambda$. Maka terdapat suatu $\alpha \in S$
sedemikian sehingga $p \in \alpha$. Karena $\alpha$ adalah suatu potongan,
terdapat suatu $q \in \alpha$ sedemikian sehingga $p < q$. Jadi $q \in
\lambda$.
\end{enumerate}
Hal ini membuktikan pernyataan tersebut. Selain itu, jelas bahwa $(\forall \alpha \in S)\alpha
\subseteq \bigcup S = \lambda$, yakni\ $\lambda$ merupakan suatu batas atas bagi $S$. Sekarang andaikan $\beta \in \mathbb{R}$ juga merupakan suatu batas atas, yakni\ $(\forall \alpha \in S)\alpha \subseteq \beta$. Untuk setiap $p \in \Rat$, jika $p \in \lambda$, maka terdapat $\alpha \in S$ sedemikian sehingga $p \in \alpha$, sehingga $p \in \beta$. Secara umum, $\lambda \subseteq \beta$. Jadi $\lambda$ merupakan batas atas \emph{terkecil} bagi $S$.
\end{proof}

Jadi, kita mempunyai sekumpulan entitas yang memenuhi Sifat Kelengkapan.
Salah satu cara untuk mengungkapkannya ialah: tidak ada ``celah'' pada
potongan-potongan kita. (Dengan demikian, membuat ``potongan'' lanjutan dari
bilangan real, bukan dari bilangan rasional, tidak akan menghasilkan objek
baru yang menarik.)

Selanjutnya, kita harus mendefinisikan beberapa operasi pada bilangan real.
Kita mulai dengan menanamkan bilangan rasional ke dalam bilangan real dengan
menetapkan bahwa $p_\Real = \Setabs{q \in \Rat}{q < p}$ untuk setiap $p \in
\Rat$. Kemudian kita mendefinisikan:
\begin{align*}
	\alpha + \beta &= \Setabs{p + q}{p \in \alpha \land q \in \beta}\\
%	\alpha - \beta &\defis \Setabs{p - q}{p \in \alpha \land q \in \Rat \setminus \beta}\\
	\alpha \times \beta &=
	\Setabs{p \times q}{0 \leq p \in \alpha \land 0 \leq q \in \beta} \cup 0_\Real & \text{jika }\alpha, \beta \geq 0_\Real
%	\alpha \div \beta &\defis
%	\Setabs{\nicefrac{p}{q}}{p \in \alpha \land q \in \Rat \setminus \beta}  & \text{jika }\alpha \geq 0_\Real, \beta > 0_\Real
\end{align*}
Untuk menangani kasus-kasus perkalian lainnya, pertama-tama tetapkan: %bahwa $0_\Real \times \alpha = 0_\Real = \alpha \times 0_\Real$, dan tambahkan:
\begin{align*}
	-\alpha &= \Setabs{p - q}{p < 0 \land q \notin \alpha}
\end{align*}
dan kemudian tetapkan:
\begin{align*}
	\alpha \times \beta &\defis
	\begin{cases}
		\mathord{-}\alpha \times \mathord{-}\beta &\text{jika }\alpha < 0_\Real\text{ dan }\beta < 0_\Real\\
		\mathord{-}(\mathord{-}\alpha \times \beta) &\text{jika }\alpha < 0_\Real \text{ dan }\beta > 0_\Real\\
		\mathord{-}(\alpha \times \mathord{-}\beta) &\text{jika }\alpha > 0_\Real \text{ dan }\beta < 0_\Real
	\end{cases}
\end{align*}
Selanjutnya, kita perlu memeriksa bahwa setiap definisi tersebut selalu
menghasilkan suatu potongan. Terakhir, kita perlu menjalani pembuktian yang
mudah (tetapi panjang lebar) bahwa potongan-potongan yang didefinisikan
demikian berperilaku tepat sebagaimana semestinya. Namun, semua ini kita
tempatkan dalam \olref[check]{sec}.
\end{document}
